\section{From prepared credit to credit on one solution}\label{sec:introduction}
The reciprocal chiral exchange extension of the Tri-Octagon model provides an explicit balance between native trial storage, a conditional periodic carrier and a reversible receiver.~\cite{baseline}
Its homogeneous native feedback matches carrier work, while receiver loading gives an equal and opposite change of receiver storage. The predecessor proved global exchange evolution, initial supply and reception from empty external stores, and native credit for other prepared carrier states.

Those results leave a distinct dynamical question: does a solution that starts with an empty carrier and receiver later credit the native store without resetting any state? A prepared negative-power state and a reversing symmetry do not answer that question for the original launch.
Here we give an explicit sufficient receiver-coupling threshold for a declared coherent family. The total native power first has the supplying sign and later has the returning sign on the same unique solution. The threshold and witness are finite; propagation and native feedback remain active in the exact equations.

The proof develops the accepted same-trajectory investigation and its internal analytical review.~\cite{investigation,review}
It uses the predecessor's global-solution and launch theorems as premises, then controls a simple comparison by an energy estimate. Section~\ref{sec:problem} specifies the unchanged candidate, Section~\ref{sec:theorem} states the result, and Sections~\ref{sec:energy}--\ref{sec:comparison} give the complete new argument.
The resulting credit is recovery of earlier supplied trial storage. Section~\ref{sec:meaning} distinguishes that fact from a recovered fraction, an attainable capacity or a spatial circuit.
